% !TeX root = ../notes.tex
% Planning version: section skeleton, content bullets, live TikZ, figure list, and reading list.

\chapter{Generating high-quality chemical data}

\begin{notebox}
  \keypoint{Intellectual question:} How do we generate the high-quality, information-rich experimental datasets that modelling and experiment selection require?
\end{notebox}

\section{More experiments are not necessarily more information}

\begin{itemize}
  \item Distinguish experiment count, measurement density, chemical diversity, and information content.
  \item Compare endpoint measurements with full reaction-progress curves.
  \item Explain that throughput can amplify systematic error just as efficiently as it amplifies useful data.
  \item Connect experimental design to the chemical question rather than to the maximum number of wells or samples.
\end{itemize}

\solution{6.5cm}{%
  \centering
  \resizebox{0.92\linewidth}{!}{%
    \begin{tikzpicture}[
      box/.style={draw=darkBlue,fill=paleBlue,line width=1.5pt,minimum width=3.8cm,minimum height=1.0cm,align=center,font=\sffamily},
      good/.style={draw=darkGreen,fill=paleGreen,line width=1.5pt,minimum width=3.8cm,minimum height=1.0cm,align=center,font=\sffamily},
      arrow/.style={->,>=Stealth,draw=darkGrey,line width=2pt},
      label/.style={font=\small\sffamily,text=darkGrey,align=center}
    ]
      \node[box] (many) at (0,1.3) {many repeated endpoints};
      \node[good] (rich) at (0,-1.3) {diverse, calibrated\\reaction profiles};
      \node[box] (count) at (6,1.3) {high throughput};
      \node[good] (info) at (6,-1.3) {high information content};
      \draw[arrow] (many) -- (count);
      \draw[arrow] (rich) -- (info);
      \node[label] at (3,-2.2) {these are related, but not equivalent};
    \end{tikzpicture}%
  }
}

\section{Practical FAIR chemical data}

\begin{itemize}
  \item Define findable, accessible, interoperable, and reusable in laboratory terms.
  \item State explicitly that FAIR does not necessarily mean open.
  \item Include identifiers, machine-readable structures, units, instrument methods, calibrations, processing steps, and provenance.
  \item Record unsuccessful, interrupted, and analytically ambiguous experiments rather than silently deleting them.
  \item Use the FAIR principles as a compact reference, not as a separate informatics lecture.\mcite{Wilkinson2016}
\end{itemize}

\solution{7.2cm}{%
  \centering
  \resizebox{0.92\linewidth}{!}{%
    \begin{tikzpicture}[
      stage/.style={draw=darkBlue,fill=paleBlue,line width=1.5pt,minimum width=3.0cm,minimum height=0.9cm,align=center,font=\sffamily},
      arrow/.style={->,>=Stealth,draw=darkGrey,line width=2pt},
      note/.style={font=\small\sffamily,text=darkGrey,align=center}
    ]
      \node[stage] (identity) at (0,2.0) {chemical identity};
      \node[stage] (conditions) at (4.2,2.0) {conditions and time};
      \node[stage] (instrument) at (8.4,2.0) {instrument method};
      \node[stage] (processing) at (8.4,0) {processing and QC};
      \node[stage] (outcome) at (4.2,0) {outcome and uncertainty};
      \node[stage] (provenance) at (0,0) {provenance and status};
      \draw[arrow] (identity) -- (conditions);
      \draw[arrow] (conditions) -- (instrument);
      \draw[arrow] (instrument) -- (processing);
      \draw[arrow] (processing) -- (outcome);
      \draw[arrow] (outcome) -- (provenance);
      \node[note] at (4.2,-1.1) {a reusable experimental record};
    \end{tikzpicture}%
  }
}

\section{High-throughput and parallel experimentation}

\begin{itemize}
  \item Introduce reaction arrays, well plates, liquid handling, and miniaturisation.
  \item Compare broad categorical screens with dense local condition screens.
  \item Include reagent identity checks, dosing accuracy, evaporation, mixing, sealing, and cross-contamination.
  \item Treat unsuccessful experiments as part of the dataset when the analytical evidence supports that label.
\end{itemize}

\section{Measurement is part of the experiment}

\begin{itemize}
  \item Show sampling, quenching, dilution, calibration, LC, GC, MS, NMR, spectroscopy, and automated processing as parts of one measurement chain.
  \item Distinguish ``no product'', ``below detection'', ``analysis failed'', and ``experiment failed''.
  \item Include blanks, standards, replicates, retention-time drift, detector saturation, and data-quality flags.
  \item Make the analytical bottleneck central to the design of automated campaigns.
\end{itemize}

\solution{5.8cm}{%
  \centering
  \resizebox{0.95\linewidth}{!}{%
    \begin{tikzpicture}[
      stage/.style={draw=darkCyan,fill=paleCyan,line width=1.5pt,minimum width=2.5cm,minimum height=0.9cm,align=center,font=\sffamily},
      arrow/.style={->,>=Stealth,draw=darkGrey,line width=2pt}
    ]
      \node[stage] (reaction) at (0,0) {reaction};
      \node[stage] (sample) at (3.4,0) {sample and\\quench};
      \node[stage] (measure) at (6.8,0) {instrument\\measurement};
      \node[stage] (process) at (10.2,0) {processing\\and calibration};
      \node[stage] (record) at (13.6,0) {QC and\\record};
      \draw[arrow] (reaction) -- (sample);
      \draw[arrow] (sample) -- (measure);
      \draw[arrow] (measure) -- (process);
      \draw[arrow] (process) -- (record);
    \end{tikzpicture}%
  }
}

\section{Flow as a data-generation platform}

\begin{itemize}
  \item Introduce only the reactor concepts needed to understand the data.
  \item Define residence time $\tau=V/Q$, steady state, lag, dispersion, and inline or at-line analysis.
  \item Explain why continuous variables are natural in flow and why data points are not independent if the system has not equilibrated.
  \item Compare parallel batch, serial flow, and hybrid architectures by task fit.
\end{itemize}

\solution{7.0cm}{%
  \centering
  \resizebox{0.94\linewidth}{!}{%
    \begin{tikzpicture}[
      box/.style={draw=darkBlue,fill=paleBlue,line width=1.5pt,minimum width=3.0cm,minimum height=1.0cm,align=center,font=\sffamily},
      analyser/.style={draw=darkGreen,fill=paleGreen,line width=1.5pt,minimum width=3.0cm,minimum height=1.0cm,align=center,font=\sffamily},
      arrow/.style={->,>=Stealth,draw=darkGrey,line width=2pt},
      label/.style={font=\small\sffamily,text=darkGrey,align=center}
    ]
      \node[box] (feed) at (0,0) {reagent feeds};
      \node[box] (reactor) at (4.3,0) {reactor\\$\tau=V/Q$};
      \node[box] (steady) at (8.6,0) {steady-state\\sampling};
      \node[analyser] (analysis) at (12.9,0) {inline or at-line\\analysis};
      \draw[arrow] (feed) -- (reactor);
      \draw[arrow] (reactor) -- (steady);
      \draw[arrow] (steady) -- (analysis);
      \node[label] at (6.4,-1.25) {lag and dispersion link neighbouring samples};
    \end{tikzpicture}%
  }
}

\section{High-throughput kinetics: Mayer}

\begin{itemize}
  \item Use the liquid-handling and continuous-UV workflow for acid/base catalysis in water.\mcite{KufferMayer2026}
  \item Emphasise full kinetic curves rather than a single endpoint.
  \item Show how a mechanistic rate model separates acid, base, buffer, and background contributions.
  \item Use the change in relative substrate reactivity across catalytic regimes as evidence that information-rich data retain mechanistic value.
  \item Bridge to Chapter~5 through leave-one-group-out validation and reactivity prediction.
\end{itemize}

\section{Automated synthesis and analysis: Bourne}

\begin{itemize}
  \item Introduce the stopped-flow architecture as an intermediate between batch and continuous flow.\mcite{Avila2022}
  \item Discuss the integration of liquid handling, reaction timing, at-line analysis, and library synthesis.
  \item Use the reported library campaign to compare throughput, reagent use, and synthesis success.
  \item Treat the feed-forward neural network as a consumer of the generated data, not as the chapter's main subject.
\end{itemize}

\section{Cage discovery: Greenaway progression}

\subsection{Computation fused with robotic synthesis}

\begin{itemize}
  \item Use the 78 precursor combinations and the complete positive, partial, and negative outcome matrix.\mcite{Greenaway2018}
  \item Compare predicted topology and formation energy with experimental outcome.
  \item Retain the unexpected cage-catenane result as evidence that systematic screening can reveal unplanned chemistry.
\end{itemize}

\subsection{Automated synthesis, characterisation, and analysis}

\begin{itemize}
  \item Use the 366-reaction campaign and the computational enumeration of hypothetical outcomes.\mcite{Basford2024}
  \item Show automated turbidity, NMR, and HRMS analysis and how orthogonal measurements assign different aspects of outcome.
  \item Use the released experiment map, code, and raw data as the strongest practical FAIR example.
  \item Discuss the reduction in manual data-analysis burden without implying that all assignments become chemically trivial.
\end{itemize}

\section{Szczypiński cage automation}

\begin{itemize}
  \item Use computation to map the accessible cage conformers.\mcite{ShieldsCage2024}
  \item Use the semi-automated high-throughput crystallisation screen to show execution at scale.
  \item Discuss how solid-state structures ``lock'' otherwise dynamic conformers and therefore create information that solution data alone do not provide.
  \item Classify this as genuine automation, but not autonomy or closed-loop optimisation.
\end{itemize}

\section{From execution to autonomy}

\begin{itemize}
  \item Close by separating machine execution from machine selection.
  \item Chapter~7 ends when the system has generated, measured, quality-controlled, and recorded data.
  \item Chapter~8 begins when those measurements determine what the platform does next.
\end{itemize}

\solution{5.8cm}{%
  \centering
  \begin{tikzpicture}[
    auto/.style={draw=darkBlue,fill=paleBlue,line width=1.6pt,minimum width=4.4cm,minimum height=1.0cm,align=center,font=\sffamily},
    autonomous/.style={draw=darkGreen,fill=paleGreen,line width=1.8pt,minimum width=4.4cm,minimum height=1.0cm,align=center,font=\sffamily},
    arrow/.style={->,>=Stealth,draw=darkGrey,line width=2pt}
  ]
    \node[auto] (execute) at (0,1.4) {automation: execute a defined plan};
    \node[auto] (record) at (0,0) {measure, QC, and record};
    \node[autonomous] (decide) at (0,-1.4) {autonomy: use the result to choose next};
    \draw[arrow] (execute) -- (record);
    \draw[arrow] (record) -- (decide);
  \end{tikzpicture}
}

\section{Summary and review}

\begin{itemize}
  \item Does the campaign increase experiment count, information content, or both?
  \item Are identities, conditions, measurements, processing, and failures reusable?
  \item Where is the analytical bottleneck?
  \item Does the platform execute a plan or choose subsequent experiments?
  \item Which parts of the workflow still require chemical judgement?
\end{itemize}

\section*{Author production checklist: publication figures}

\begin{itemize}
  \item Kuffer and Mayer: Fig.~3 for the high-throughput kinetic workflow; Fig.~4 for mechanism-guided decomposition; Fig.~6 only as the modelling bridge.
  \item Avila \textit{et al.}: Fig.~1 for the stopped-flow platform; retain one later figure showing library success and one showing model validation, with numbering checked against the downloaded publisher PDF before extraction.
  \item Greenaway \textit{et al.} (2018): Fig.~3 for precursor space, computation, and the complete experimental outcome matrix.
  \item Basford \textit{et al.} (2024): Fig.~3 for the integrated experimental/computational workflow; Fig.~5 for all 366 automated outcomes; Fig.~6 for the outcome Sankey summary.
  \item Shields \textit{et al.} (2024): Fig.~2 for the computed conformer landscape; Fig.~3 for the five experimentally observed conformations.
\end{itemize}

\section*{Author reading and PDF-download list}

\begin{itemize}
  \item FAIR guiding principles.\mcite{Wilkinson2016}
  \item High-throughput information-rich kinetics.\mcite{KufferMayer2026}
  \item Automated stopped-flow library synthesis.\mcite{Avila2022}
  \item Computation fused with robotic cage synthesis.\mcite{Greenaway2018}
  \item Automated porous-organic-cage discovery and analysis.\mcite{Basford2024}
  \item Semi-automated cage crystallisation and conformer discovery.\mcite{ShieldsCage2024}
\end{itemize}
